% !TEX program = pdflatex
\documentclass[11pt]{article}
\usepackage{support-notes}
\title{Bounded profiles and finite-dimensional points of universal localisations}
\author{\workingauthor}\date{\coursedate}
\begin{document}\maketitle
\begin{abstract}
We classify every bounded profile of a positive graded Koszul dual
having a suitable normal element, a hereditary free algebra on the open
chart, and perfect lifts of its finite simple points. One part of the
order consists of finite orbit sets of simple modules. The other is given
by the finite-dimensional points of universal localisations at orbit
matrices. Their mixed order changes at free rank one. All hypotheses
are verified for arbitrary periods, numbers of generators and finite-order
linear twists, over arbitrary fields. Statements labelled as theorems,
lemmas, corollaries and examples are proved.
\end{abstract}

\section{Assumptions and statement}
Use cochain grading. Let $A$ be a connected locally finite
nonnegatively graded $k$-algebra with zero differential, and
$h\in A$ a homogeneous regular normal element of degree $n>0$.
Put $G=A/(h)$ and $B=A[h^{-1}]$. Normality and regularity give a
degree-preserving automorphism $\theta$ with $ha=\theta(a)h$ and a
two-sided Ore localization flat on each side (as either module it is
a direct limit of shifted free rank-one modules). For a perfect module
call its cohomology \emph{eventually invertible under $h$} if right or
left multiplication as appropriate gives bijections in all sufficiently
high degrees. We use the following hypotheses.
\begin{enumerate}
\item $G$ has infinite total dimension. It admits a degree-zero outer
derivation $G\to G\otimes k\ell\otimes G$ taking a positive homogeneous
element $v$ to $1\otimes\ell\otimes1$, $|v|=|\ell|$.
Alternatively $h$ is even and central and on the fibre a tensor of two
perfect factors is finite only if one factor has finite cohomology.
\item $B\cong R[u,u^{-1};\sigma]$ as a graded algebra with $|u|=1$,
$R=B^0=k\langle U_1,\ldots,U_g\rangle$, $g\ge1$, and $uf=\sigma(f)u$.
For the poset classification suppose $\sigma$ has finite order $L$.
\item Each finite-dimensional simple right or left $R$-module is the
degree-zero localized cohomology of a perfect $A$-module on the matching
side with eventually invertible $h$.
\end{enumerate}
There is no restriction on the ground field. The derivation in (1) is a
checkable sufficient condition; only the implication in Lemma~\ref{detect}
is used in the classification. Hypothesis (3) ensures surjectivity of
the proposed labels and supplies sources distinguishing them.

Denote by $\Omega_r,\Omega_l$ the finite-dimensional simple types of
$R$ on the indicated side modulo twisting by $\sigma$. For a finite
module $W$, let $S(W)$ be the set of orbits of its composition factors.
For a square matrix $H$ over $R$ put
\[
 O(H)=\{[T]\in\Omega_l:\ H(T^{\sigma^a})\text{ invertible for all }
 0\le a<L\}.
\]
Left evaluation means the usual block action on columns of elements of
$T$; on the right use precomposition, hence rows. Write $P^\vee=\RH_A(P,A)$
for the left perfect dual of a right perfect source. Profiles are the
subcategories of $\per A$ on which $\RH_A(P,E)$ has finite total
cohomology. All such Hom complexes have locally finite, bounded-below
cohomology by finite cells and summands.

\begin{thm}\label{main}
All profiles occur exactly once in the following two rows, with equality
of the displayed subset as label equivalence:
\[
 \begin{array}{c|c|c}
  &\text{labels}&\text{inclusion}\\ \hline
 \mathcal T_F& F\subset\Omega_r\text{ finite} &\text{reverse inclusion of }F\\
 \mathcal B_O& O=O(H)\subseteq\Omega_l &\text{ordinary inclusion of }O .
 \end{array}
\]
The first row consists of detectors $E$ with eventually invertible $h$,
labelled by $S(H^0E_h)$. Sources in $\mathcal B_O$ are exactly those for
which $h$ is eventually invertible on $P^\vee$ and $S(H^0(P^\vee)_h)
\subseteq O$. The top is $\mathcal T_\varnothing$ and bottom is
$\mathcal B_\varnothing$. No first-row profile lies in a second-row profile.
For $g\ge2$, a second-row profile lies in a first-row profile exactly
when the former is the bottom or the latter is the top. For $g=1$
identify left and right simple orbits via the closed points of
$\mathbb A^1_k$. Then $\mathcal B_O\subseteq\mathcal T_F$ if and only if
$O\cap F=\varnothing$.
\end{thm}

\section{The closed and open tests}
\begin{lem}\label{detect}
For perfect right $E$ and perfect left $Y$, finite cohomology of
$E\ten_A Y$ is equivalent to both
\[
 \begin{gathered}
 h\text{ eventually invertible on }H^*E\text{ or on }H^*Y,\\
 H^{i}(E_h\ten_B Y_h)=0\quad(i\gg0).
 \end{gathered}
\]
This lemma requires only hypothesis (1) and the normal localization.
\end{lem}
\begin{proof}
If neither outer module satisfies the first condition, the relative
two-cut estimate of \othernote{two_cuts}{A lower bound for derived tensor
products} supplies $\dim H^{b+j}(E\ten_A Y)\ge\dim G^j$ for fixed
$b$ and all $j\ge0$. It applies because $A\to G$ is onto, the chosen
element lifts to a cycle, and regularity identifies reduction images
with cokernels of $h$. Under the alternate central hypothesis the
same implication follows from Theorem 1 of
\othernote{central_reduction}{Central reduction for bounded profiles}.
For one parameter, finite generation over $k[h]$ is equivalent to
eventual bijectivity: beyond the generator degrees multiplication is
surjective, and the finite dimensions then stabilize in each residue
class. Conversely lift bases of the finitely many low degrees.
If $E$ satisfies the condition, $E\to E_h$ is a high-degree cohomology
isomorphism, since the localization of a graded module at a regular normal
element on the ring is computed by fractions. Tensor its cone with $Y$.
Vanishing in high degrees is preserved by the finite cells and retracts
in $Y$. Associativity identifies the second tensor with
$E_h\ten_B Y_h$. The same argument with the sides reversed applies if
$Y$ satisfies the condition. This proves both implications.
\end{proof}

For finite-order $\sigma$ the homogeneous unit $u^{2L}$ is even and
graded central in $B$. It makes the localized tensor periodic, turning
high-degree vanishing into acyclicity. This uses no centrality for an
odd $h$. Even without periodicity, if the right factor $E$ is eventually
invertible, one may evaluate $\RH_B(P_h,E_h)$ in high degrees.
If the left factor $P^\vee$ is eventually invertible, one may instead
evaluate $E_h\ten_B(P^\vee)_h$ there. Perfect duality gives these
identifications on both charts.

\begin{lem}\label{chart}
Every perfect dg module over $B$ is quasi-isomorphic to its graded
cohomology with zero differential and to the cone of a finite matrix
of degree-zero free modules $B^a\to B^b$. If the cohomology of the
original object over $A$ has eventually invertible $h$, then the
degree-zero part of its localized cohomology is finite dimensional.
\end{lem}
\begin{proof}
Taking degree zero is an equivalence of graded right $B$-modules with
right $R$-modules, inverse $W\mapsto W\otimes_R B$; multiplication
by $u$ proves the componentwise assertion. Consecutive components are
twists, and there is an analogous left equivalence. Submodules of free
$R$-modules are free. For a finite ambient basis, choose leading words
using a right-compatible degree-lexicographic order with a module
position. Choose one monic vector for each prefix-minimal leading word
in its position. Their right multiples are independent by distinct
leading words and span by leading-word reduction. In particular the
image of every finite matrix is finite free, the kernel splits off
the domain and is finite projective, hence free. This gives coherence.
The left statement follows by reversing words. It follows by the long
exact sequence and coherence (kernels, cokernels and extensions) that
cohomology of any perfect dg $B$-module is finitely presented as a
graded module. Retracts preserve finite presentation.

Present this cohomology by an injective matrix $B^a\to B^b$ over $R$,
using the equivalence to place every generator in degree zero. Lift the
free generators to cycles and each relation to a primitive of its image.
The resulting map from the cone induces an isomorphism on cohomology.
The cone also maps quasi-isomorphically to the graded cokernel. Finally
localization agrees with eventually invertible cohomology in high
degrees, which are finite dimensional, and $u$ identifies all components.
\end{proof}

\section{Proof of the classification}
For an eventually invertible detector, apply Hom from the presentation
of $P_h$ into the graded module induced by $W=H^0(E_h)$. The map between
free terms in each degree is evaluation of the matrix on a twist of $W$
(up to cochain signs which are units). Its cone is acyclic exactly when
the map is bijective in every twist. A rectangular matrix of unequal
ranks cannot be bijective on a nonzero finite module. For square
matrices a composition series gives a block triangular evaluation with
square diagonal blocks. Hence evaluation on all simple factors and their
twists is necessary and sufficient. This proves one implication of
reverse inclusion. For the other, consider an orbit present in $F'$
and absent from $F$. The annihilator of a finite-dimensional simple
is maximal two-sided: its faithful image is finite-dimensional primitive,
hence simple artinian. Equal annihilators identify the simple on the
same side. The Chinese remainder theorem for the finitely many distinct
maximal two-sided ideals of the chosen types and all twists gives
$f\in R$ acting by one on the types of $F$ and by zero on the chosen
orbit. No absolute simplicity or separability was used. Clear the
denominator by multiplying on the left by $h^d$, giving an element
of $A^{dn}$. The cone of this left multiplication between shifted
free right $A$-modules supplies a witness. On localization the extra
homogeneous unit only shifts the set of twists being tested. Direct
sums of lifts in (3) realise all finite labels, including the empty one.

For every other detector Lemma~\ref{detect} requires $h$ eventually
invertible on $Y=P^\vee$. Write $T=H^0Y_h$. Use
Lemma~\ref{chart} to present $E_h$ instead and tensor with the graded
left module induced by $T$. A rectangular unequal-size presentation
forces $T=0$, and a square presentation $H$ gives precisely
$S(T)\subseteq O(H)$. All labels are realised: clear the denominators of
$H$ with a common $h^d$ and take as detector $G$ plus the cone of the
resulting matrix over $A$. The $G$ summand has zero localization and
infinite cohomology killed by $h$; as a right or left $A$-module it
is perfect by regularity. This construction also works for a matrix
with kernel. The identity and zero matrices give $\Omega_l$ and the
empty domain. Duals of left lifts in (3) distinguish any two unequal
domains. Perfect double duality follows on cells and summands.
The source $G$ belongs to every first-row profile and no second-row
profile: its right resolution is a single free multiplication cone,
so Hom from it tests eventual invertibility of $h$ on the detector.

To check mixed inclusions, suppose $P$ lies in the second row and test
it against a first-row detector with finite module $W$. A finite
$d$-dimensional right $R$-module $W$ has the free presentation
\[
 0\longrightarrow W\otimes k^g\otimes R\longrightarrow W\otimes R
 \longrightarrow W\longrightarrow0,qquad
 w\otimes i\otimes r\longmapsto w\otimes U_i r-wU_i\otimes r .
\]
Leading first letters prove injectivity and word reduction proves
exactness in the middle. On a nonzero left $T$, the dimensions of its
two free terms after evaluation are $gd\dim T$ and $d\dim T$. For
$g\ge2$, $d>0$, acyclicity is impossible. The intersection is therefore
exactly the sources with $T=0$ and eventually invertible $h$, namely
$\mathcal B_\varnothing=\UU_A$. In rank one use the primary decomposition
over the PID $k[U]$. The resolution by a power of a monic irreducible
has zero derived tensor with a different primary module and nonzero
cokernel with a module for the same prime. Running over all twists
shows that acyclicity is equivalent to disjoint orbit supports. All
finite supports can be tested by the lifts. This proves
Theorem~\ref{main}. For $g=1$, nonempty domains are precisely the
complements of finite orbit sets, by determinants over the PID. With
identity twist, labelling the second row by the excluded points plus
infinity identifies the whole order with finite subsets of the closed
points of $\mathbb P^1_k$ in reverse inclusion, together with a new
bottom. Over a finite field there are still infinitely many closed points.

\section{Universal localisation and its precise limitation}
For a square matrix $H$ on free right $R$-modules let $R_{\Sigma(H)}$
be Cohn's universal localisation, with
$\Sigma(H)=\{\sigma^a H:0\le a<L\}$. A finite-dimensional left module,
viewed as a representation into its $k$-endomorphism ring, extends to
the localisation precisely when each block evaluation is invertible.
Indeed the inverse block operators are forced, and adjoining their
entries subject to both inverse identities constructs the extension.
Thus the class of all such finite modules has exactly those finite
simple factors in $O(H)$, by the square triangular test. It follows
that the second row is the poset of essential images on
finite-dimensional left modules of these orbitwise square-matrix
localisations, ordered by inclusion. A finite collection of matrices
is combined by a block sum. The assertion concerns these images;
\href{https://arxiv.org/abs/0810.1401}{Krause--\v Stov\'i\v cek}
identify universal localisations with homological epimorphisms for
hereditary rings, a result about the full module categories.

\begin{ex}[Two different localisations with the same profile label]
Take two disjoint cyclic generator orbits $u_a,v_a$ of equal length.
Invert all $u_a v_a$, or invert all $u_a$ and $v_a$ individually.
On finite-dimensional left modules these conditions coincide. On a
countably infinite-dimensional space let every $u_a$ be the backward
shift and every $v_a$ the forward shift. The products act as identity,
but neither shift is invertible. This representation extends only to
the first localisation. The induced ring epimorphisms are inequivalent
under $R$, including after orbit closure. This excludes any unrestricted
bijection with the epimorphisms themselves.
\end{ex}

\begin{ex}[A hereditary open chart for which extensions matter]
Let $R$ be the path algebra of $1\to2$, put $A=R[t]$ with $|t|=2$,
and use left modules in this example. The representations
$(k\xrightarrow{1}k)$ and $(k\xrightarrow{0}k)$ have the same
composition factors. Extend them by $k[t]$; both are perfect since $R$
is hereditary. Hom from the cone of the map of vertex projectives
representing the arrow evaluates that arrow in every positive even
degree. One detector gives finite Hom and the other gives infinite
Hom. Maps between unlike projectives need not induce square diagonal
blocks on simple factors. Heredity alone cannot justify the first-row
simple-factor test.
\end{ex}

\section{A class with arbitrary periods and numbers of generators}
For $n,r\ge1$ and $Q\in\mathrm{GL}_r(k)$ consider
\[
 A_{n,r,Q}=k\langle x,z_1,\ldots,z_r\rangle/
 \big(x^n z_i-\textstyle\sum_j Q_{ij}z_j x^n\big),\quad
 |x|=|z_i|=1 .
\]
\begin{cor}\label{family}
If $Q$ has finite order, Theorem~\ref{main} applies with $h=x^n$,
$g=nr$ and the twists of the free generators described below. For
arbitrary $Q$, Lemmas~\ref{detect} and~\ref{chart} and the perfect lifts
still hold with the indicated infinitely many possible twists.
These rings occur as residue-field Koszul duals of proper augmented
dg algebras. For $n=2$ the nonpositive kernels can be taken two-term.
\end{cor}
\begin{proof}
Rewrite $x^n z_i$ to $\sum_j Q_{ij}z_jx^n$. Length-lex order with $x$
largest decreases. No proper overlap among leading words is possible,
as each ends in a $z$ and all its nontrivial proper prefixes are powers
of $x$. Disjoint replacements commute linearly; induction on order
gives a unique normal form. Normal words consist of blocks $x^a z_i$,
$0\le a<n$, followed by a power of $x$. Right multiplication by $x^n$
is injective. Moving $x^n$ past letters implements the invertible
linear substitution $Q$ and fixes $x$, giving normality and left
regularity. The quotient by $x^n$ is $(k[x]/x^n)*k\langle z_1,\ldots,z_r
\rangle$, with the required free partial derivative and nonzero
$z_1$-powers in every degree.

In the localization put $U_{i,a}=x^a z_i x^{-(a+1)}$.
Conjugation by $x$ moves $a$ to $a+1$ and sends the last column to
$\sum_j Q_{ij}U_{j,0}$. These rules define an automorphism of the free
algebra on $nr$ generators. The skew Laurent algebra so defined maps
to the localization by the displayed formulas. In the other direction
send $z_i$ to $U_{i,0}x$ and $x$ to the Laurent generator; the maps
respect the relations and are inverse. If $Q^d=1$ then $\sigma^{nd}=1$.

Given any finite right module $W$, in the induced localized module
take its nonnegative-degree part $M$. It is generated by its degree-zero
basis. A presentation has these generators and, for each basis vector
$w$, the $nr$ relations
\[
         w x^a z_i-(wU_{i,a})x^{a+1},\qquad 0\le a<n .
\]
Their right multiples are independent: filter normal words by the number
of $z$'s, use the maximal input filtration degree, and observe that the
leading outputs begin with distinct initial blocks and a specified basis
vector. All corrections have one fewer $z$, also after reduction. The
series of $A$ is $1/((1-t)(1-r\sum_{j=1}^n t^j))$, by the word basis.
Thus the graded cokernel has series $\dim(W)/(1-t)$ and maps onto
$M$, which has exactly that series. This is a finite free resolution.
The opposite algebra has the reversed relation with inverse linear
substitution; the same argument gives left lifts. Multiplication by
$x$ on these nonnegative tails is bijective between successive
components in their allowed degrees.

For completeness take the strictly unital augmented minimal
$A_\infty$-algebra whose reduced basis consists of $r+1$ degree-zero
letters and $r$ outputs of degree $1-n$. Give its reduced bar dual the
free generators $x,z_i$ of degree one and $y_i$ of degree $n$ with
$d y_i=x^nz_i-\sum_jQ_{ij}z_jx^n$, $dx=dz_i=0$.
This specifies the signs and identities: the graded derivation squares
to zero; equivalently the only reduced operation has arity $n+1$ and
its outputs never re-enter it. For $n=1$ this specifies ordinary
radical multiplication. At fixed weight (letters weight one, $y_i$
weight $n+1$) filter by the lexicographic expansion $y_i\mapsto x^n z_i$.
The leading differential expands only these monomials. In a fixed
expanded word its disjoint occurrences each supply a contractible
two-term choice, with tensor signs from homogeneous degrees. Words
with no occurrences survive in degree equal to weight and are already
cycles. The finite filtration proves that the bar dual projects
quasi-isomorphically to $A_{n,r,Q}$. Degreewise bar duals are ordinary
here because generator degrees are positive.

The bar--cobar strictification of the finite minimal model is
nonpositive: a cobar generator on a nonempty bar word of $b$ letters
has degree the sum of their cohomological degrees minus $b-1$.
It can even be chosen finite-dimensional. Use the internal weights
above on the original letters and outputs. The bar--cobar comparison
to the minimal model is weight preserving on cohomology, so the
subcomplex in weights $>n+1$ is acyclic and is a two-sided ideal.
Truncating it leaves finitely many words and the same quasi-isomorphism
type. The comparison is the usual augmented bar--cobar equivalence,
as in \href{https://arxiv.org/abs/math/9910179}{Keller's treatment of
minimal models and strictification}. The finite-cohomology category is
generated by the residue field by truncation and finite length of the
local degree-zero heart. When $n=2$, a concrete strict two-term family
is given in the companion notes (for general $Q$ use $b_i r=
\sum_j Q_{ji} w_j$, $rb_i=w_i$ in the strict table there).
\end{proof}

For $Q=I,n=2,r=1$ this reproduces the two rows, the swap and the mixed
order from the earlier six-dimensional example. For larger $n$ it is a
theorem about proper dg algebras; a two-term kernel realisation by an
ordinary universal localisation is claimed only when the cohomological
amplitude satisfies the manuscript's realisation theorem. The
factorisation of complex free loci by
\href{https://arxiv.org/abs/1907.04328}{Helton--Klep--Vol\v ci\v c}
further labels nonempty domains by finite sets of stable atomic factors,
saturated under the finite twist, as in the original example. The
arbitrary-field classification above uses no algebraic closure.

\begin{ex}[A symmetric disjoint-support rule is impossible]
In this family with $nr\ge2$, let $M_0$ lift the one-dimensional right
zero tuple and take $G=A/(x^n)$ as another right perfect. Then
$\RH_A(G,M_0)$ is finite and $\RH_A(M_0,G)$ is infinite. For the
first claim the resolution of $G$ tests $x^n$, which is bijective on
the tail of $M_0$. For the second, the displayed resolution of $M_0$
gives in degree $i$ a map from $G^i$ to
$\bigoplus_{a=0}^{n-1}\bigoplus_{j=1}^r G^{i+a+1}$.
Appending $z_1$ embeds the normal words of $G^i$ in $G^{i+1}$,
and $z_1^i$ is nonzero. The target is strictly larger in every
nonnegative degree and the cokernels survive in the Hom cone.
Thus no assignment of a single set to each object can detect this
bounded Hom relation by disjointness of the two assigned sets.
\end{ex}
\end{document}
